% 来源：sources/2016-2025-期中-甲-历年真题汇编.pdf 第 11 页
\begin{problem}{}{6cm}
计算行列式
$D_1=\begin{vmatrix}2&4&6&8\\4&6&8&2\\6&8&2&4\\8&2&4&6\end{vmatrix}$。
\end{problem}

\begin{problem}{}{9cm}
已知 $n+1$ 阶行列式
\[
\setlength{\arraycolsep}{4.5pt}
D=\begin{vmatrix}
-1&x_1&x_2&\cdots&x_n\\
1&2&0&\cdots&0\\
1&0&3&\cdots&0\\
\vdots&\vdots&\vdots&&\vdots\\
1&0&0&\cdots&n+1
\end{vmatrix},
\]
$x_1,x_2,\cdots,x_n$ 为任意常数，求
\[
4A_1+9A_2+\cdots+(n+1)^{2}A_n,
\]
其中 $A_1,A_2,\cdots,A_n$ 分别为 $x_1,x_2,\cdots,x_n$ 的代数余子式。
\end{problem}

\begin{problem}{}{8.5cm}
求解线性方程组
\[
\begin{cases}
\lambda x_1+x_2+x_3+x_4=1,\\
x_1+\lambda x_2+x_3+x_4=\lambda,\\
x_1+x_2+\lambda x_3+x_4=\lambda^{2}.
\end{cases}
\]
\end{problem}

\begin{problem}{}{7.5cm}
用初等行变换将矩阵
$A=\begin{pmatrix}1&2&5&-1&-1\\1&2&3&1&3\\3&6&13&-1&1\\5&10&21&-1&3\end{pmatrix}$
化简成行阶梯形矩阵 $B$，使 $B$ 的阶梯头为\mbox{~$1$}。
\end{problem}

\begin{problem}{}{9.5cm}
设 $r$ 为自然数。
\begin{enumerate}[label=(\arabic*),leftmargin=2.2em,itemsep=0.25em,topsep=0.3em]
  \item 叙述秩的定义；
  \item 证明：$r(A)\le r$ 当且仅当 $A$ 中 $r+1$ 阶子式（若有）均为 $0$；
  \item 设 $A$ 通过一次行倍加得到 $B$，求证：$r(B)=r(A)$。
\end{enumerate}
\end{problem}

\begin{problem}{}{11cm}
设有 $n\ (\ge 2)$ 阶方阵 $A$，$A^{*}$ 为 $A$ 的伴随矩阵，证明：
\begin{enumerate}[label=(\arabic*),leftmargin=2.2em,itemsep=0.3em,topsep=0.3em]
  \item $r(A^{*})=\begin{cases}n,&r(A)=n,\\1,&r(A)=n-1,\\0,&r(A)<n-1;\end{cases}$
  \item $(A^{*})^{*}=|A|^{n-2}A$；
  \item 若 $r(A)>0$，则 $(A^{*})^{*}=A$ 当且仅当 $|A|^{n-2}=1$。
\end{enumerate}
\end{problem}

\begin{problem}{}{6.5cm}
设 $A,C,D$ 均为 $m\times n$ 阶矩阵，$B,D$ 均为 $n\times s$ 阶矩阵，证明：$r(AB-CD)\le r(A-C)+r(B-D)$。
\end{problem}

\begin{problem}{}{8.5cm}
设 $A$ 为 $2022^{2022}$ 阶方阵，用数学归纳法证明：
\[
\forall n\in Z^{+},\ \forall m\in Z^{+},\quad (n+1)r(A^{2})\le r(A^{n+2})+mr(A).
\]
\end{problem}
